% Generated draft: requires mathematical and source review.
\documentclass[11pt]{article}
\usepackage[margin=25mm]{geometry}
\usepackage{fontspec}
\setmainfont{Latin Modern Roman}
\usepackage{amsmath,amssymb,mathtools,unicode-math}
\setmathfont{Latin Modern Math}
\usepackage{enumitem,needspace,xurl,hyperref}
\setlength{\parindent}{0pt}
\setlength{\parskip}{5pt}
\setlength{\emergencystretch}{3em}
\newcommand{\N}{\mathbb{N}}
\newcommand{\Z}{\mathbb{Z}}
\newcommand{\Q}{\mathbb{Q}}
\newcommand{\R}{\mathbb{R}}
\newcommand{\C}{\mathbb{C}}
\newcommand{\F}{\mathbb{F}}
\newcommand{\A}{\mathbb{A}}
\newcommand{\PP}{\mathbb{P}}
\newcommand{\E}{\mathbb{E}}
\newcommand{\eps}{\varepsilon}
\newcommand{\dd}{\,\mathrm{d}}
\begin{document}
\section*{UnsolvedMath ID: 700398}
\textbf{Source problem code:} catalogue entry
\subsection{Definitions and mathematical statement}
Let \(p\) be a prime, and let \(G=F\rtimes A\) be a Frobenius group with Frobenius kernel \(F\) and complement \(A=\langle a\rangle\) of order \(p\). Thus \(F\trianglelefteq G\), \(G=FA\), \(F\cap A=\{1\}\), and the complement acts on \(F\) without nonidentity fixed points: \(C_F(b)=\{1\}\) for every \(b\in A\setminus\{1\}\), where \(C_F(b)=\{f\in F:fb=bf\}\). Write \(a^g=g^{-1}ag\).

A group \(H\) is \emph{locally finite} if \(\langle X\rangle\) is finite for every finite subset \(X\subseteq H\). Interpret \emph{binary finite} as meaning that \(\langle x,y\rangle\) is finite for every \(x,y\in H\). The supplied record does not define this term, so this interpretation requires confirmation against the cited source.

The problem asks whether each of the following implications holds for every such Frobenius group, with the two parts considered separately:
\begin{enumerate}
\item[(a)] If \(G\) is binary finite, must \(G\) be locally finite? Equivalently, does
\[
\bigl(\forall x,y\in G,\ |\langle x,y\rangle|<\infty\bigr)
\quad\Longrightarrow\quad
\bigl(\forall X\subseteq G\text{ finite},\ |\langle X\rangle|<\infty\bigr)
\]
hold?
\item[(b)] If \(\langle a,a^g\rangle\) is finite for every \(g\in G\), must \(F\) be locally finite? That is, does
\[
\bigl(\forall g\in G,\ |\langle a,a^g\rangle|<\infty\bigr)
\quad\Longrightarrow\quad
\bigl(\forall X\subseteq F\text{ finite},\ |\langle X\rangle|<\infty\bigr)
\]
hold? No binary-finiteness assumption on \(G\) is imposed in this part.
\end{enumerate}
\subsection{Short English statement}
For a Frobenius group with a cyclic complement of prime order, ask separately whether finiteness of every two-generated subgroup forces the whole group to be locally finite, and whether finiteness of every subgroup generated by the complement generator and one of its conjugates forces the kernel to be locally finite. The first question uses the stated interpretation of ``binary finite.''
\subsection{Sources}
\begin{itemize}
\item Supplied local UnsolvedMath registry record: permanent ID \(700398\), exact source problem code \texttt{catalogue entry}, title \emph{Kourovka Notebook Problem 6.56}.
\item Attributed to V. P. Shunkov, proposed in 1978. The supplied source is \emph{The Kourovka Notebook (Unsolved Problems in Group Theory)}, 6th issue (1978), Problem 6.56, as recorded in the 21st issue edition (2026), \texttt{arXiv:1401.0300v48}. The collection description names E. I. Khukhro and V. D. Mazurov as editors.
\item The registry labels the problem \emph{open}; its supplied update timestamp is \texttt{2026-10-16T14:00:00Z}. This is supplied registry information, not an independently verified assessment of the literature or current mathematical status.
\end{itemize}
\end{document}
